\documentclass[10pt,a4paper]{amsart}
\usepackage[margin=19mm]{geometry}
\usepackage{amsmath,amssymb,mathtools}
\usepackage[scr=rsfs,cal=euler]{mathalfa}
\usepackage{xfrac}
\usepackage[unicode,hidelinks,bookmarks=false]{hyperref}
\newcommand{\ie}{i.e.\ }
\newcommand{\cf}{cf.\ }
\newcommand{\resp}{respectively\ }
\newcommand{\Subsection}{\subsection}
\providecommand{\nobreakdash}{\nobreak\hbox{-}}
\setlength{\emergencystretch}{2em}
\multlinegap0pt
\usepackage{xcolor}
\usepackage[T1]{fontenc}
\setcounter{MaxMatrixCols}{20}
\newcommand{\Z}{\mathbb{Z}}
\newcommand{\blank}{\rule{.3cm}{0.15mm}}
\newcommand{\dslash}{/\mkern-5mu/}
\newcommand{\dbslash}{\backslash\mkern-5mu\backslash}
\DeclareMathOperator{\Hom}{Hom}
\newcommand{\Mlt}{\mathop{\mathrm{Mlt}}}
\newcommand{\Sym}{\mathop{\mathrm{Sym}}}
\newcommand{\Aut}{\mathop{\mathrm{Aut}}}
\newcommand{\Inn}{\mathop{\mathrm{Inn}}}
\newtheorem{theorem}{Theorem}[section]
\newtheorem{lemma}[theorem]{Lemma}
\newtheorem{definition}[theorem]{Definition}
\newtheorem{example}[theorem]{Example}
\newtheorem{proposition}[theorem]{Proposition}
\newtheorem{claim}[theorem]{Claim}
\newtheorem{corollary}[theorem]{Corollary}
\newtheorem{conjecture}[theorem]{Conjecture}
\newtheorem{problem}[theorem]{Problem}
\newtheorem{observation}[theorem]{Observation}
\newtheorem{counterexample}[theorem]{Counterexample}
\newtheorem{exercise}[theorem]{Exercise}
\newtheorem{note}{Note}[theorem]
\newtheorem{fact}{Fact}[theorem]
\newtheorem{comment}[theorem]{Comment}
\newtheorem{question}[theorem]{Question}
\newtheorem{List}[theorem]{List}
\theoremstyle{definition}
\newtheorem{remark}[theorem]{Remark}
\title[Homology of quasigroups of Bol-Moufang type]{In search of homology for quasigroups of Bol-Moufang type: the boundary identity $(\partial_2 \circ Q)(T) = - \bar{g}(r) + \sum_{l \in L(T)} t^{i(l)} s^{j(l)}\bar{g}(l)$ (Lemma 4.1)}
\author{Anthony Christiana, Ben Clingenpeel, Huizheng Guo, Jinseok Oh, J\'{o}zef H. Przytycki and Anna Zamojska-Dzienio}
\date{Source: arXiv:2508.21268v1 (CC BY 4.0)}
\begin{document}
\begin{abstract}
We initiate (co)homology theory for quasigroups of Bol-Moufang type based on analysis of their extensions by affine quasigroups of the same type. We use these extensions to define second and third boundary operations, $\partial_2(x,y)$ and $\partial_3(x,y,z)$, respectively. We use these definitions to compute the second homology groups for several examples from the work of Phillips and Vojtechovsky. We speculate about the relation between these homology groups and those obtained from a small category with coefficients in a functor.
\end{abstract}
\maketitle

\noindent\emph{Source and licence.} In search of homology for quasigroups of Bol-Moufang type. Version arXiv:2508.21268v1 (CC BY 4.0), distributed under the Creative Commons Attribution 4.0 International licence, http://creativecommons.org/licenses/by/4.0/, as recorded in the arXiv OAI licence metadata for this exact version and shown on the abstract page https://arxiv.org/abs/2508.21268v1. Full rights notice: licenses/arxiv-2508.21268-CC-BY-4.0.txt.\par\smallskip
\noindent\textit{Editorial scope.} Counted unit: the paper's boundary identity for the tree polynomial $Q$, printed as Lemma 4.1 in the published version (source0.tex lines 2059--2062, unlabelled in the source), delivered together with the paper's complete original proof of it (source0.tex lines 2064--2090) as the single primary proof body. No mathematics is written, repaired or re-derived: statement and proof are the paper's own text line for line, in the paper's own notation ($Q$, $g$, $\bar g$, $\partial_2$, $i(v)$, $j(v)$, $V_L$, $v_{LL}$, $v_{LR}$). Required context retained uncounted, in document order: Section 2.3 and its Definition of an affine quasigroup together with the remark recording that medial quasigroups are exactly the affine quasigroups with commuting automorphisms (source0.tex lines 379--398); the lemma collecting the affine identities $(a+b)*c=a*c+tb=b*c+ta$, $a*(b+c)=a*b+sc=a*c+sb$ and $(a+b)*(c+d)=a*c+b*d-c_0$ (source0.tex lines 613--621); the subsubsection on equivalent extensions with the definition of equivalence of two extensions, the derivation of $\partial_2(x,y)=tx+sy-xy$ and Proposition 3.4 with its proof (source0.tex lines 661--698); and Section 4 with its subsection on the chain-complex verification, which recalls $Q(T)=\sum_{v\in V(T)-L(T)} h_v(t,s)g(v)$, the recursive definition of $\bar g(v)$ and its companion $g(v)=(\bar g(v_L),\bar g(v_R))$, and the meaning of $i(v)$ and $j(v)$ (source0.tex lines 2027--2056). Every definition, proposition, remark, numbered display, footnote and citation of those lines is retained; the two paragraphs that follow the counted proof in the source (the statement of Corollary 4.2 at source0.tex lines 2094--2098 and the discussion at lines 2100--2102) are not part of this unit and are omitted, so the delivered document ends with the counted proof. Not part of this unit and not redistributed: the front matter (title, author list with addresses, emails, keywords and subject classification, source0.tex lines 60--86, of which the title and the author list are reproduced in the wrapper title block and the abstract is reproduced verbatim in the wrapper front matter); the table of contents; Sections 1 and 2.1--2.2 (source0.tex lines 95--378); the subsubsections on the extension calculus for right Bol quasigroups and on substitutions and rooted binary trees (source0.tex lines 599--660 and 706--2058 except the four retained spans); the whole of the Examples section, the speculation section, the acknowledgments and the appendix tables (source0.tex lines 2103--3638); and the paper's bibliography listing, of which only the cited entries are transcribed. The delivered text cites four keys (Shch; and Bruc, Mur and Toy) through exactly the three citation commands that the retained lines contain: a citation of Shch carrying the section locator 1.4.4.3, a joint citation of Bruc, Mur and Toy, and a citation of Shch carrying the locator Theorem 2.65. Those four entries are transcribed verbatim from the paper's own in-source thebibliography (source0.tex lines 3363--3486) with the paper's printed bibliography labels preserved (Shch, Bruc, Mur, Toy), so every printed citation label of the delivered unit is the published one. Reference adaptation: the retained context line 2034 refers to Definition 3.5 of the paper, which lies outside every retained span. That single cross-reference is delivered as a hyperlink to the pinned version with the published number kept literally, so the delivered sentence reads ``Recall from Definition 3.5 the polynomial $Q$ on a binary rooted tree $T$'' with the number 3.5 hyperlinked; this is the only reference adaptation made. All other labels and references inside the delivered unit resolve inside it: the label of the retained affine-quasigroup definition, the label lem:32 of the retained identities, the label Second of the retained Proposition 3.4 together with the commented-out paragraph that refers to it, the label partialsquared of the retained subsection 4.1, and the paper's own section, subsection and theorem counters, which are set so that the printed numbers coincide with the published version: Definition 2.2 for the affine quasigroup, Remark 2.3 for the medial remark, Lemma 3.2 for the affine identities, Definition 3.3 and Proposition 3.4 for the equivalent-extension block, subsection 4.1 for the chain-complex verification, and Lemma 4.1 for the counted unit. One printed heading is not claimed to coincide with the published version: the retained subsubsection heading on equivalent extensions is set with the wrapper's own subsubsection counter. Printed slips kept literal: no printed word, symbol, hypothesis or number of the counted statement or of its proof was altered. The paper's own forms are delivered as printed, among them the use of $\bar g(v)\in X$ and $g(v)\in X^2$ for elements that the paper elsewhere treats as group-ring elements, its unlabelled lemma environment, its habit of writing the product of two non-leaf children as $\bar g(v_L)\bar g(v_R)$ in the definition and as $\bar g(v_L)*\bar g(v_R)$ in the proof, its coloured blue and red terms used as an aid to the reader (which is why xcolor is loaded), its mixed use of the symbol $\subset$ for both proper and improper inclusion (declared in a footnote on the retained context), and its own spelling of the reference keys. Non-ASCII characters: the paper contains five characters outside ASCII in one distinct code point, all five of them en dashes inside the entries Eil, EM, EM12 and Leb-2 of its bibliography (source0.tex lines 3385, 3392, 3406 and 3434); none of those entries is cited by the retained text and none of them is transcribed, so the delivered unit contains no character outside ASCII, and no encoding substitution, no glyph replacement and no missing-character risk arises. Layout: the paper's own class is amsart at 11pt (source0.tex line 1) with graphicx, multicol, makecell, chngpage, tabularx, xcolor, amsmath, amssymb, amsthm, mathtools, tikz-cd and hyperref; the wrapper is the lane's pinned amsart editorial wrapper at 10pt on A4, which already supplies amsmath, amssymb, mathtools and hyperref, and to which this package adds xcolor and the paper's own macro block (source0.tex lines 35--44: the MaxMatrixCols setting and the paper's nine custom macros, which define the symbol for the integers, a blank rule, the two slash symbols and the operators Hom, Mlt, Sym, Aut and Inn) and the paper's own theorem declarations (source0.tex lines 11--32) verbatim, so every printed number is produced by the paper's own counters. multicol, makecell, chngpage, tabularx and tikz-cd are not loaded because no retained span uses them. No publisher class, no publisher style file and no upstream script is redistributed: the paper is already an amsart article and the wrapper is editorial formatting. Assets: the paper has five figure environments and no graphics-inclusion command anywhere (its five figures are drawn with the picture environment and with a tabular); no figure environment lies inside any retained span, no retained span contains a graphics-inclusion command, no retained span contains a PDF-page inclusion command, and no retained span contains a TeX input or include directive; the package declares no asset dependency, and no image file is copied into or redistributed with this package. All mathematics of the unit is therefore native text with no image dependency. A rights notice is delivered at licenses/arxiv-2508.21268-CC-BY-4.0.txt.
\setcounter{section}{2}\setcounter{subsection}{2}\setcounter{subsubsection}{0}\setcounter{theorem}{1}
\subsection{Affine quasigroups}\label{Ss:AMBMQ} %of Bol-Moufang quasigroups}
In the next sections, we will consider extensions of Bol-Moufang quasigroups by affine quasigroups of the same type. For this purpose we have to know the structure of affine BMq's. 
\begin{definition}
Let $(A,+)$ be an abelian group, $t,s\colon A\rightarrow A$ (not necessarily commuting) automorphisms of $(A,+)$ and $c_0\in A$. We define a new operation on $A$ by
\begin{equation}
a\ast b=t(a)+s(b)+c_0.    
\end{equation}
The resulting quasigroup $(A,\ast)$ is said to be \emph{affine} over $(A,+)$.
\end{definition}
Affine quasigroups are also called \emph{central quasigroups} or \emph{$T$-quasigroups} (\cite[Section 1.4.4.3]{Shch}). They are polynomially equivalent to a module over the ring of Laurent polynomials of two non-commuting variables.

If $(A,\ast)$ is an affine quasigroup, then the left division is defined by $a\backslash b=s^{-1}\big(b-ta-c_0\big)$ and the right division by $a/ b= t^{-1}\big(a-sb-c_0\big)$. 

\begin{remark}
A quasigroup $(Q,\cdot)$ is \emph{medial} (or \emph{entropic}) if it satisfies the medial law
\begin{equation}
(x\cdot y)\cdot (z\cdot t)=(x\cdot z)\cdot (y\cdot t).
\end{equation}
The famous Bruck-Murdoch-Toyoda Theorem \cite{Bruc,Mur,Toy} (see also \cite[Theorem 2.65]{Shch}) states that, up to isomorphism, medial quasigroups are precisely affine quasigroups with commuting automorphisms $t,s$.
\end{remark}

\setcounter{section}{3}\setcounter{subsection}{1}\setcounter{subsubsection}{0}\setcounter{theorem}{1}
We used in our calculation the following simple lemma:
\begin{lemma}\label{lem:32}
We have the following identities in the affine quasigroup $(A,\ast)$ over an abelian group $(A,+)$:
\begin{enumerate}
\item[(1)] $(a+b)*c= ta+tb+sc+c_0= ta+sc+c_0 + tb= a*c+tb= b*c+ ta$.
\item[(2)] $ a*(b+c)= ta+sb +sc+ c_0= a*b+sc= a*c+sb$.
\item[(3)] $ (a+b)*(c+d)= ta+tb +sc+ sd+ c_0= a*c+tb+sd=a*c+b*d-c_0$.
\end{enumerate}
\end{lemma}

\setcounter{theorem}{2}\setcounter{subsubsection}{2}
\subsubsection{Equivalent extensions and $\partial_2$} 

We consider in this subsection the conditions under  which two extensions by an affine quasigroups are equivalent. This will allow us to recover $\partial_2(x,y): kX^2\to kX$.
\begin{definition} 
We say that two extensions of $X$ by $A$, where $(A\times X,*_1)$ is given by a $2$-cochain $\phi_1$ and $(A\times X,*_2)$ is given by a $2$-cochain  $\phi_2$ are equivalent if there is a $1$-cochain $\alpha: X \to A$ such that the map 
$\hat\alpha: (A\times X,*_1)\to (A\times X,*_2)$ given by 
$$\hat\alpha(a,x)= (a+\alpha(x),x)$$
is a quasigroup homomorphism. That is:
$$\hat\alpha((a_1,x_1)*_1(a_2,x_2))=\hat\alpha(a_1,x_1)*_2\hat\alpha(a_2,x_2)$$
\end{definition}
We will show that the condition for a homomorphism (in fact an isomorphism) $\hat\alpha$ leads to the second boundary map $\partial_2$ given by:
$$\partial_2(x,y)= tx + sy - xy.$$
See Section~\ref{partialsquared} for a proof that $\partial_2 \partial_3 = 0$. 
%\textcolor{red}{Note this is where we're defining $\partial_2$, we should maybe say a little more (but not too much more) and/or point readers to this proposition from future sections dealing with $\partial_2$.}

Recall that in $A$ we have $a*b= ta+sb+c_0.$

\begin{proposition}\label{Second}
Two equivalent extensions $(A\times X,*_1)$ and  $(A\times X,*_2)$ given by $2$-cochains $\phi_1$ and $\phi_2$ and related using isomorphism $\hat\alpha$ lead to the equation:
$$\phi_1(x_1,x_2)- \phi_2(x_1,x_2)=\alpha(tx_1+sx_2-x_1x_2).$$

\end{proposition}
\begin{proof}
We compute
$$\hat\alpha((a_1,x_1)*_1(a_2,x_2))= \hat\alpha(a_1*a_2+\phi_1(x_1,x_2),x_1x_2)=$$
$$(a_1*a_2+ \alpha(x_1x_2)+ \phi_1(x_1,x_2),x_1x_2)$$

$$\hat\alpha(a_1,x_1)*_2\hat\alpha(a_2,x_2)=(a_1+\alpha(x_1),x_1)*_2 (a_2+\alpha(x_2),x_2)= $$
$$((a_1+\alpha(x_1))*_2(a_2+\alpha(x_2))+ \phi_2(x_1,x_2),x_1x_2)\stackrel{Lemma \ref{lem:32}}{=}$$
$$(a_1*a_2+ t\alpha(x_1)+s\alpha(x_2)+ \phi_2(x_1,x_2),x_1x_2).$$
For these two expressions to be equal we need:
$$\phi_1(x_1,x_2)-\phi_2(x_1,x_2)= t\alpha(x_1) +s\alpha(x_2)- \alpha(x_1x_2).$$
Equivalently we can write it as:
$$\phi_1(x_1,x_2)-\phi_2(x_1,x_2)= \alpha(tx_1+sx_2- x_1x_2)$$ 
and from this we deduce that
$\partial_2(x_1,x_2)= tx_1+sx_2-x_1x_2$.

\end{proof}

\setcounter{section}{3}\setcounter{subsection}{0}\setcounter{subsubsection}{0}\setcounter{theorem}{0}
\section{BMq Chain Complexes and Their Properties}\label{S:chain}

\subsection{Chain complex verification}\label{partialsquared}
In this subsection, we verify that 
\[ 0 \to C_3 \xrightarrow{\partial_3} C_2 \xrightarrow{\partial_2} C_1 \to 0 \]
 is, in fact, a chain complex. We perform the verification of $\partial_2 \partial_3 = 0$ in slightly greater generality than our situation demands.

 Recall from Definition~\href{https://arxiv.org/abs/2508.21268v1}{3.5} the polynomial $Q$ on a binary rooted tree $T$:



\[ Q(T) = \sum_{v \in V(T) - L(T)} h_v(t,s) g(v) \]

where $g(v) = (g(v_L), g(v_R))$ with 
\[
\bar g(v)=
\left\{
\begin{array}{ll}
x_i & \mbox{when $v\in L(T)$ is decorated by $x_i$},\\
\bar g(v_L)\bar g(v_R) & \mbox{when $v\in V(T)-L(T)$ and $v_L$, $v_R$} \\
& \mbox{are the left/right children of $v$ in $T$}.
\end{array}
\right.
\]
 
This map is a generalization of $\partial_3: C_3 \to C_2$ in that 
$\partial_3(x,y,z) = Q(T_1) - Q(T_2)$ for two trees $T_1, T_2$ with $4$ leaves decorated by the same 4-letter word in $x,y,z$. 

Finally, recall that $i(v), j(v)$ are the number of left (resp. right) turns from the root vertex $r$ to the vertex $v.$ 


\begin{lemma}
    For  $r \in V$ the root vertex,
    \[ (\partial_2 \circ Q)(T) = - \bar{g}(r) + \sum_{l \in L(T)} t^{i(l)} s^{j(l)}\bar{g}(l) \]
\end{lemma}

\begin{proof}
    For the sake of clarity, we note that $\bar g (v) \in X$, $g(v) \in X^2$, and $\bar g = x_i$ for some $x_i \in W$.


By definition, 

\begin{align*}
    (\partial_2 \circ Q)(T) & = \sum_{v \in V(T)-L(T)} t^{i(v)}s^{j(v)} \partial_2(g(v)) \\
    & = \sum_{v \in V(T)-L(T)} t^{i(v)}s^{j(v)} (t \bar g(v_L) + s \bar g (v_R) - \bar g(v_L) * \bar g(v_R)) \\
     & = \sum_{v \in V(T)-L(T)} t^{i(v) +1}s^{j(v)} \bar g(v_L) + t^{i(v)}s^{j(v)+1} \bar g (v_R) - t^{i(v)}s^{j(v)}  \bar g(v) \\
\end{align*}

Assume for $v \in V(T)- L(T)$ that $v_L \in V_L$; i.e. the left child of $v$ is also a degree three vertex.

Then 
\[ \partial_2(g(v)) = {\color{blue} t^{i(v) +1}s^{j(v)} \bar g(v_L)} + {\color{red} t^{i(v)}s^{j(v)+1} \bar g (v_R) } - t^{i(v)}s^{j(v)}  \bar g(v)\]

 and 
 
\[ \partial_2(g(v_L)) = t^{i(v) +2}s^{j(v)} \bar g(v_{LL}) + t^{i(v)+1}s^{j(v)+1} \bar g (v_{LR}) {\color{blue} - t^{i(v)+1}s^{j(v)}  \bar g(v_L) } \]
 have a pair of cancelling terms in blue. Similarly the term in red will cancel with the final term in $\partial_2(g(v_R))$ so long as $v_R \in V(T)- L(T)$. 


If a vertex $v \in V(T)-L(T)$ has a child which is not leaf, say $g(v_L) = x_i$, we produce a non-cancelling term  $t^{i(v_L)} s^{j(v_L)} \bar g(v_L)$. 

Finally, we note that the final term in $\partial_2(g(r))$ (namely, $-\bar{g}(r)$), will not cancel. Thus we arrive at the statement.
\end{proof}

\begin{thebibliography}{99}
\bibitem[Shch]{Shch} V. Shcherbacov, Elements of quasigroup theory and applications,
CRC Press, 2017.
\bibitem[Bruc]{Bruc} R.~H.~Bruck, Some results in the theory of quasigroups, {\it Trans. Amer. Math. Soc.} vol. 55, 1944, pp. 19-52
\bibitem[Mur]{Mur} D.~C.~Murdoch,
Quasi-groups which satisfy certain generalized laws, {\it American
Journal of Math.}, 16, 1939, pp. 509-522.
\bibitem[Toy]{Toy} K.~Toyoda,
On affine geometry of abelian groups,
{\it Proceedings of the Imperial Academy}, Volume 16(5),  1940, 161-164.
\end{thebibliography}
\end{document}
